\title{Group Sequence Policy Optimization}

\begin{abstract}
We present Group Sequence Policy Optimization (GSPO), a reinforcement learning algorithm tailored for large language model training, offering improvements in stability, efficiency, and performance. In contrast to previous approaches that operate using token-level importance ratios, GSPO leverages a sequence-level perspective—calculating importance ratios across entire sequences, and applying clipping, reward assignment, and optimization at this level. Our empirical results indicate that GSPO not only surpasses the GRPO algorithm in both sample efficiency and model performance, but also enhances the stability of Mixture-of-Experts (MoE) reinforcement learning setups and streamlines RL system design. The adoption of GSPO has played a pivotal role in the significant advancements observed in the latest Qwen3 models.
\end{abstract}

\section{Introduction}
\label{sec:intro}

Reinforcement learning (RL) has become a foundational approach for enhancing and expanding the capabilities of language models \citep{o1,dpsk-r1,qwq32b,qwen3}. Leveraging large-scale RL enables these models to address complex tasks, including advanced mathematics and programming at a competitive level, by engaging in extended and more profound reasoning.

Effective scaling of RL, particularly with increased computational resources, requires that training remains stable and resilient. Despite this necessity, contemporary leading RL algorithms—such as GRPO \citep{grpo}—are prone to pronounced instability during the training of massive language models, frequently culminating in catastrophic and irreversible failure modes \citep{qwen3, minimax-m1}. Such instability significantly impedes the ongoing advancement of language model capabilities through RL.

In this work, we demonstrate that GRPO’s instability primarily results from an incorrect and invalid utilization of importance sampling weights within the algorithm, which causes significant variance in training noise. This variance accumulates as response sequences grow longer and is exacerbated by the use of clipping, ultimately leading to model degradation.

To overcome these foundational challenges, we introduce \textbf{Group Sequence Policy Optimization (GSPO)}, a novel RL framework designed for large language model training. GSPO’s central contribution is its rigorous and principled computation of the importance ratio using sequence likelihood, as established in prior literature \citep{click}, thereby adhering closely to the original purpose of importance sampling. Further, GSPO calculates normalized rewards as the comparative advantages among multiple generated responses to a single query, promoting coherence between reward assignment and sequence-level optimization.

Through empirical analysis, we establish that GSPO greatly outperforms GRPO in terms of training robustness, speed, and overall efficacy. Notably, GSPO intrinsically addresses the instability issues commonly observed in RL training for large-scale Mixture-of-Experts (MoE) architectures, obviating the reliance on elaborate stabilization methods and indicating opportunities for streamlining RL systems. The advantages afforded by GSPO were instrumental in realizing marked performance gains in the recent Qwen3 models. Looking ahead, we regard GSPO as a dependable and scalable cornerstone capable of fostering future progress in large language model RL training.

\section{Preliminaries}

\paragraph{Notation}
Throughout this work, we represent an autoregressive language model with parameters $\theta$ as a policy $\pi_\theta$. The variable $x$ refers to a query, and $\mathcal{D}$ indicates the collection of queries. For a response $y$ generated for query $x$, its probability under the policy $\pi_\theta$ is expressed as $\pi_\theta (y | x)=\prod_{t=1}^{|y|} \pi_\theta (y_t | x, y_{<t} )$, where $|y|$ corresponds to the total token count in $y$. Each query-response tuple $(x, y)$ is evaluated by a verifier $r$, producing a reward $r(x, y)$ that lies in the interval $[ 0, 1 ]$.

\paragraph{Proximal Policy Optimization (PPO)} Using samples generated from the old policy $\pi_{\theta_\text{old}}$, PPO \citep{ppo} constrains the policy update within a proximal region of the old policy through the clipping mechanism. Specifically, PPO employs the following objective for policy optimization (we omit the KL regularization term hereinafter for brevity, as it is not the focus of this paper):

\begin{align}
\mathcal{J}_\text{PPO}(\theta) = \mathbb{E}_{ x \sim \mathcal{D},\, y \sim \pi_{\theta_\text{old}}( \cdot | x) }
\left[ \frac{1}{|y|} \sum_{t=1}^{|y|} 
\min \left( w_{t}(\theta) \widehat{A}_{t},  \, \mathrm{clip} \left( w_{t}(\theta), 1 - {\varepsilon}, 1 + {\varepsilon}\right) \widehat{A}_{t} \right)
\right],
\end{align}

In this formulation, the importance ratio for a token $y_t$ is given by $
w_{t}(\theta) = \frac{ \pi_{\theta} (y_{t} | x, y_{<t}) }{ \pi_{\theta_\text{old}} (y_{t} | x,y_{<t})}
$, with the advantage $\widehat{A}_{t}$ of $y_t$ computed using an auxiliary value model, and the parameter $\varepsilon$ specifying the permissible range for the ratio's clipping.

A significant obstacle encountered when applying PPO arises from its dependency on the value model. Typically, this value model mirrors the policy model in scale, contributing substantially to both computational and memory demands. In addition, the algorithm’s performance is closely linked to the accuracy of its value predictions. Developing a trustworthy value model is a complex endeavor, and maintaining its efficacy as tasks increase in complexity or sequence length constitutes an even more difficult problem.

\paragraph{Group Relative Policy Optimization (GRPO)} GRPO \citep{grpo} bypasses the need for the value model by computing the relative advantage of each response within a group of responses to the same query. Specifically, GRPO optimizes the following objective:

\begin{align}
\label{equ:grpo}
\mathcal{J}_\text{GRPO}(\theta) = \mathbb{E}_{ x \sim \mathcal{D},\, \{y_i\}_{i=1}^G \sim \pi_{\theta_\text{old}}( \cdot | x) }
\left[ \frac{1}{G} \sum_{i=1}^{G} \frac{1}{|y_i|} \sum_{t=1}^{|y_i|} 
\min \left( w_{i,t}(\theta) \widehat{A}_{i,t},  \, \mathrm{clip} \left( w_{i,t}(\theta), 1 - {\varepsilon}, 1 + {\varepsilon}\right) \widehat{A}_{i,t} \right)
\right],
\end{align}

where $G$ is the number of generated responses to each query $x$ (i.e., the group size), and the importance ratio $w_{i,t}(\theta)$ and advantage $\widehat{A}_{i,t}$ of token $y_{i,t}$ are:

\begin{align}
    w_{i,t}(\theta)=\frac{ \pi_{\theta} (y_{i,t} | x, y_{i,<t}) }{ \pi_{\theta_\text{old}} (y_{i,t} | x,y_{i,<t})},\quad\ 
    \widehat{A}_{i,t} = \widehat{A}_{i} = \frac{r(x, y_i) - \mathrm{mean} \left( \{ r(x, y_i) \}_{i=1}^G \right) }{ \mathrm{std} \left( \{ r(x, y_i) \}_{i=1}^G \right) },
\end{align}

respectively, where all the tokens in $y_i$ share the same advantage as $\widehat{A}_{i}$.

\section{Motivation}
\label{sec:motivation}

As model architectures increase in size and complexity—such as through the use of sparsity techniques like Mixture-of-Experts models—and as generated response sequences lengthen, it becomes essential to utilize large rollout batch sizes for efficient hardware performance in reinforcement learning. In order to make training more sample-efficient, the typical approach involves dividing these large batches of rollout data into several smaller mini-batches, which are then used for performing gradient updates. However, this procedure gives rise to an off-policy learning context, since the response samples $y$ are drawn from a previous policy $\pi_{\theta_\text{old}}$ rather than the present policy $\pi_\theta$ that is undergoing optimization. This scenario highlights why mechanisms such as clipping in PPO and GRPO are crucial; they limit the influence of samples that deviate significantly from the current policy, thereby ensuring that gradient estimates are not adversely affected by excessively off-policy data.

While mechanisms like clipping aim to manage this off-policy discrepancy, we identify a more fundamental issue in GRPO: \textit{its objective is ill-posed}. This problem becomes particularly acute when training large models on long-response tasks, leading to catastrophic model collapse. The ill-posed nature of the GRPO objective stems from a misapplication of importance sampling weights. The principle of importance sampling is to estimate the expectation of a function $f$ under a target distribution $\pi_\text{tar}$ by re-weighting samples drawn from a behavior distribution $\pi_\text{beh}$:

\begin{align}
\mathbb{E}_{z \sim \pi_\text{tar}} \left[ f(z) \right] = \mathbb{E}_{z \sim \pi_\text{beh}} \left[ \frac{ \pi_\text{tar} (z) }{ \pi_\text{beh} (z)} \, f(z) \right].
\end{align}

A key point here is that the effectiveness of importance weighting hinges on taking an average over many samples ($N \gg 1$) drawn from the behavior distribution $\pi_\text{beh}$, allowing the weight $\frac{ \pi_\text{tar} (z) }{ \pi_\text{beh} (z)}$ to adequately compensate for discrepancies between distributions.

By contrast, GRPO incorporates the importance weight $\frac{ \pi_{\theta} (y_{i,t} | x, y_{i,<t}) }{ \pi_{\theta_\text{old}} (y_{i,t} | x,y_{i,<t})}$ at each individual token step $t$. This calculation uses only a solitary sample $y_{i,t}$ from the next-token distribution $\pi_{\theta_\text{old}}( \cdot | x, y_{i,<t})$ at every position. As a result, the intended correction for the distributional shift does not materialize. Instead, this approach injects significant high-variance noise into the computed training gradients, with this noise compounding across extended sequences—an effect made more severe by the use of clipping. Our empirical findings indicate this can precipitate model collapse, which often cannot be undone. Even after reverting to prior checkpoints, diligently adjusting hyperparameters (such as clipping boundaries), increasing sequence length, or changing the RL queries, continued training fails to restore model performance after collapse.

These results indicate a deeper design limitation within GRPO. The observed failure of token-level importance weighting highlights an essential insight: \textbf{the granularity of the optimization objective ought to correspond with the granularity of the reward signal}. As rewards are associated with entire sequences, enforcing off-policy corrections on a per-token basis introduces problematic misalignment. This realization prompts us to move away from token-level objectives in favor of deploying importance weights—and carrying out the optimization process—on the \textit{sequence level} directly.

\section{Algorithm}

\subsection{GSPO: Group Sequence Policy Optimization}

Although the token-level importance weight $\frac{ \pi_{\theta} (y_{i,t} | x, y_{i,<t}) }{ \pi_{\theta_\text{old}} (y_{i,t} | x,y_{i,<t})}$ poses issues in GRPO, we note that, for language generation tasks, the \textit{sequence-level} importance weight $\frac{ \pi_{\theta} (y | x) }{ \pi_{\theta_\text{old}} (y | x)}$ carries a well-defined theoretical interpretation. Specifically, it quantifies the extent to which a response $y$ drawn from $\pi_{\theta_\text{old}} (\cdot | x)$ diverges from the distribution $\pi_{\theta} (\cdot | x)$, thereby providing a sequence-level metric that not only resonates with the overall reward structure but also informs the behavior of the clipping procedure.

Based on this straightforward observation, we propose the \textbf{Group Sequence Policy Optimization (GSPO)} algorithm. GSPO employs the following sequence-level optimization objective:

\begin{align}
\mathcal{J}_\text{GSPO} (\theta) =
\mathbb{E}_{ x \sim \mathcal{D},\, \{y_i\}_{i=1}^G \sim \pi_{\theta_\text{old}}( \cdot | x) }
\left[ 
\frac{1}{G} \sum_{i=1}^{G}
\min \left( s_{i}(\theta)  \widehat{A}_{i},  \, \mathrm{clip} \left( s_{i}(\theta), 1 - {\varepsilon}, 1 + {\varepsilon} \right) \widehat{A}_{i} \right) 
\right],
\label{equ:gspo}
\end{align}

where we adopt the group-based advantage estimation:

\begin{align}
\widehat{A}_{i} = \frac{r(x, y_i) - \mathrm{mean} \left( \{ r(x, y_i) \}_{i=1}^G \right) }{ \mathrm{std} \left( \{ r(x, y_i) \}_{i=1}^G \right) },
\end{align}

and define the importance ratio $s_{i}(\theta)$ based on sequence likelihood \citep{click}:

\begin{align}
s_{i}(\theta) = \left( \frac{ \pi_{\theta} (y_i | x) }{ \pi_{\theta_\text{old}} (y_i | x)} \right)^{\frac{1}{|y_i|}}
=
\exp \left( \frac{1}{|y_i|} \sum_{t=1}^{|y_i|} \log \frac{ \pi_{\theta} (y_{i,t} | x, y_{i,<t}) }{ \pi_{\theta_\text{old}} (y_{i,t} | x,y_{i,<t})} \right).
\end{align}

Consequently, GSPO implements clipping across whole responses, rather than at the token level, to prevent excessively ``off-policy'' samples from influencing the gradient calculation. This approach aligns with both the sequence-level reward structure and the optimization objective. Additionally, we incorporate length normalization in $s_{i}(\theta)$, aiming to stabilize variance and maintain $s_{i}(\theta)$ within a consistent numerical interval. Without length normalization, slight modifications in the probability of several tokens could trigger significant volatility in the sequence-level importance ratio, necessitating different clipping thresholds for responses of varying lengths. Furthermore, it should be recognized that the clipping intervals used in GSPO and other prior methods (such as GRPO) usually differ by orders of magnitude, owing to their respective formulations of importance ratios.

\subsection{Gradient Analysis}
\label{subsec:gradient}

We can derive the gradient of the GSPO objective as follows (clipping is omitted for brevity):

\begin{align}
\nabla_{\theta} \mathcal{J}_\text{GSPO} (\theta)
=&\ 
\nabla_{\theta} \mathbb{E}_{ x \sim \mathcal{D},\, \{y_i\}_{i=1}^G \sim \pi_{\theta_\text{old}}( \cdot | x) }
\left[ 
\frac{1}{G} \sum_{i=1}^{G} s_{i}(\theta) \widehat{A}_{i}
\right] \\
= &\ 
\mathbb{E}_{ x \sim \mathcal{D},\, \{y_i\}_{i=1}^G \sim \pi_{\theta_\text{old}}( \cdot | x) }
\left[ 
\frac{1}{G} \sum_{i=1}^{G}
s_{i}(\theta) \widehat{A}_{i}
\cdot \nabla_{\theta} \log s_{i}(\theta)
\right] \\
=&\ 
\mathbb{E}_{ x \sim \mathcal{D},\, \{y_i\}_{i=1}^G \sim \pi_{\theta_\text{old}}( \cdot | x) }
\left[ 
\frac{1}{G} \sum_{i=1}^{G} 
\left( \frac{ \pi_{\theta} (y_i | x) }{ \pi_{\theta_\text{old}} (y_i | x)} \right)^{\frac{1}{|y_i|}} \widehat{A}_{i} 
\cdot \frac{1}{|y_i|} \sum_{t=1}^{|y_i|} \nabla_{\theta} \log \pi_{\theta} (y_{i,t} | x, y_{i,<t}) 
\right].
\label{equ:gspo_grad}
\end{align}

For comparison, the gradient of the GRPO objective is as follows (note that $\widehat{A}_{i,t} = \widehat{A}_{i}$):

\begin{align}
\nabla_{\theta} \mathcal{J}_\text{GRPO} (\theta) 
=&\ 
\nabla_{\theta} \mathbb{E}_{ x \sim \mathcal{D},\, \{y_i\}_{i=1}^G \sim \pi_{\theta_\text{old}}( \cdot | x) }
\left[ \frac{1}{G} \sum_{i=1}^{G} \frac{1}{|y_i|} \sum_{t=1}^{|y_i|} 
w_{i,t}(\theta) \widehat{A}_{i,t}
\right] \\
=&\
\mathbb{E}_{ x \sim \mathcal{D},\, \{y_i\}_{i=1}^G \sim \pi_{\theta_\text{old}}( \cdot | x) }
\left[ \frac{1}{G} \sum_{i=1}^{G} \widehat{A}_{i} 
\cdot \frac{1}{|y_i|} \sum_{t=1}^{|y_i|} 
\frac{ \pi_{\theta} (y_{i,t} | x, y_{i,<t}) }{ \pi_{\theta_\text{old}} (y_{i,t} | x,y_{i,<t})} 
\nabla_{\theta} \log \pi_{\theta} (y_{i,t} | x, y_{i,<t})  
\right].
\end{align}

The primary difference between GSPO and GRPO centers on \textit{the manner in which gradients of token log likelihoods are weighted}. While GRPO assigns weights to tokens based on the ``importance weight'' $\frac{ \pi_{\theta} (y_{i,t} | x, y_{i,<t}) }{ \pi_{\theta_\text{old}} (y_{i,t} | x,y_{i,<t})}$, these weights are not uniform and may take values in $(0, 1+\varepsilon]$ when $\widehat{A}_{i} >0$ or $[1-\varepsilon, +\infty)$ when $\widehat{A}_{i} < 0$. These variable weights, far from being trivial, can accumulate throughout training, potentially introducing unpredictable instability. In contrast, GSPO ensures each token within a generated response receives the same weight, thereby mitigating this source of instability present in GRPO.

\subsection{GSPO-token: A Token-level Objective Variant}

In scenarios like multi-turn RL, we may desire a finer-grained advantage adjustment than the sequence level. To this end, we introduce a token-level objective variant of GSPO, namely \textbf{GSPO-token}, to allow token-wise advantage customization:

\begin{align}
\mathcal{J}_\text{GSPO-token}(\theta)
=
\mathbb{E}_{ x \sim \mathcal{D},\, \{y_i\}_{i=1}^G \sim \pi_{\theta_\text{old}}( \cdot | x) }
\left[ 
\frac{1}{G} \sum_{i=1}^{G} \frac{1}{|y_i|} \sum_{t=1}^{|y_i|} 
\min \left( s_{i,t}(\theta) \widehat{A}_{i,t},  \, \mathrm{clip} \left( s_{i,t}(\theta), 1 - {\varepsilon}, 1 + {\varepsilon} \right) \widehat{A}_{i,t} \right) 
\right],
\label{equ:gspo-token}
\end{align}

where

\begin{align}
s_{i,t}(\theta) 
= \mathrm{sg} \left[ s_{i}(\theta) \right]  \cdot \frac{ \pi_{\theta} (y_{i,t} | x, y_{i,<t}) }{ \mathrm{sg} \left[ \pi_{\theta} (y_{i,t} | x, y_{i,<t}) \right] },
\end{align}

and $\mathrm{sg}[\cdot]$ denotes only taking the numerical value but stopping the gradient, corresponding to the \texttt{detach} operation in PyTorch. The gradient of GSPO-token can be derived as:

\begin{align}
\nabla_{\theta} \mathcal{J}_\text{GSPO-token}(\theta)
=&\ 
\nabla_{\theta} \mathbb{E}_{ x \sim \mathcal{D},\, \{y_i\}_{i=1}^G \sim \pi_{\theta_\text{old}}( \cdot | x) }
\left[ 
\frac{1}{G} \sum_{i=1}^{G}
\frac{1}{|y_i|} \sum_{t=1}^{|y_i|} 
s_{i,t}(\theta) \widehat{A}_{i,t}
\right] \\
=&\ 
\mathbb{E}_{ x \sim \mathcal{D},\, \{y_i\}_{i=1}^G \sim \pi_{\theta_\text{old}}( \cdot | x) }
\left[ 
\frac{1}{G} \sum_{i=1}^{G} s_i (\theta) \cdot \frac{1}{|y_i|} \sum_{t=1}^{|y_i|} 
\widehat{A}_{i,t} \frac{ \nabla_{\theta} \pi_{\theta} (y_{i,t} | x, y_{i,<t}) }{ \pi_{\theta} (y_{i,t} | x, y_{i,<t}) }
\right] \\
=&\ 
\mathbb{E}_{ x \sim \mathcal{D},\, \{y_i\}_{i=1}^G \sim \pi_{\theta_\text{old}}( \cdot | x) }
\left[ 
\frac{1}{G} \sum_{i=1}^{G}  \left( \frac{ \pi_{\theta} (y_i | x) }{ \pi_{\theta_\text{old}} (y_i | x)} \right)^{\frac{1}{|y_i|}}
\cdot \frac{1}{|y_i|} \sum_{t=1}^{|y_i|}
\widehat{A}_{i,t} \nabla_{\theta} \log \pi_{\theta} (y_{i,t} | x, y_{i,<t})  
\right].
\label{equ:gspo-token_grad}
\end{align}

It is important to observe that the expression $\frac{ \pi_{\theta} (y_{i,t} | x, y_{i,<t}) }{ \mathrm{sg} \left[ \pi_{\theta} (y_{i,t} | x, y_{i,<t}) \right] }$ evaluates to 1, implying that $s_{i,t}(\theta)$ is numerically equivalent to $s_{i}(\theta)$.
Upon examining Equation~\eqref{equ:gspo} and~\eqref{equ:gspo-token}, as well as Equation~\eqref{equ:gspo_grad} alongside~\eqref{equ:gspo-token_grad}, it becomes evident that GSPO-token yields identical numerical results to GSPO in terms of the optimization objective, the clipping constraint, and the theoretical gradient provided that the advantage is assigned uniformly to every token within the response $y_i$ (specifically, $\widehat{A}_{i,t} = \widehat{A}_{i}$). However, GSPO-token extends greater adaptability, as it allows the advantage value to be tailored for each individual token.

\section{Experiments and Discussion}

\subsection{Empirical Results}

We conduct experiments using a cold-start model fine-tuned from Qwen3-30B-A3B-Base, presenting both the reward trajectories during training and the model’s performance on the AIME’24 (measured as average Pass@1 across 32 samples), LiveCodeBench (from 202410 to 202502, assessed by average Pass@1 over 8 samples), and CodeForces (using Elo Rating) evaluation suites. Throughout the RL training procedure, each batch of rollouts is subdivided into four separate mini-batches to compute gradient updates. For the GSPO algorithm, we apply left and right clipping thresholds of 3e-4 and 4e-4, as specified in Equation~\eqref{equ:gspo}. In our comparative analysis, we utilize GRPO as the baseline approach, setting its respective clipping limits in Equation~\eqref{equ:grpo} to 0.2 and 0.27; these hyperparameters have been meticulously optimized to enable an equitable comparison. It is important to highlight that GRPO relies on the Routing Replay training method to achieve stable MoE RL convergence—this aspect is further examined in \S~\ref{subsec:routing_replay}—whereas \textit{GSPO eliminates the dependence on this strategy}.

As illustrated in Figure~\ref{fig:results}, the GSPO training process maintains stability from start to finish. Our findings indicate that \textbf{GSPO consistently enhances model performance by increasing the training compute, periodically refreshing the query set, and allowing longer generation sequences}. Additionally, GSPO exhibits greater efficiency compared to GRPO, producing superior accuracy and benchmark outcomes given equivalent training compute and query usage. Finally, we have effectively utilized GSPO for RL training of the newest Qwen3 models, robustly validating GSPO’s capacity to harness the scalability benefits of RL in large language model development.

\subsection{Curious Observation on Clipping Fractions}

One fundamental difference between GSPO and GRPO lies in GSPO's method of clipping responses at the sequence level rather than operating on individual tokens. Figure~\ref{fig:clipping} illustrates this distinction, highlighting that GSPO clips tokens at a rate nearly two orders of magnitude greater than GRPO; notably, modifications to the clipping ranges do not affect this pronounced disparity. Nonetheless, even though GSPO excludes a substantially higher proportion of tokens from the training process (or from gradient estimation), it demonstrates improved training efficiency over GRPO. This surprising outcome—whereby more extensive token clipping actually enhances training efficiency—suggests that the token-level gradient calculations employed by GRPO are intrinsically noisy and suboptimal for effective sample exploitation. In contrast, the sequence-centric strategy in GSPO delivers a more dependable and potent learning signal.

\subsection{Benefit of GSPO for MoE Training}
\label{subsec:routing_replay}

\paragraph{Background}
Relative to the reinforcement learning training of dense models, MoE models exhibit distinctive challenges regarding training stability, mainly due to their sparse activation mechanism. Notably, our investigations reveal that, under the application of the GRPO algorithm, MoE models are susceptible to \textit{expert-activation volatility}, which can undermine RL training convergence. Specifically, we observe that the subset of experts engaged for a given output can shift considerably following one or more gradient updates. For instance, when training the 48-layer Qwen3-30B-A3B-Base, approximately 10\% of the experts selected for the identical rollout sample after a gradient step are not the same as those under the preceding policy $\pi_{\theta_\text{old}}$, when evaluated under the updated policy $\pi_{\theta}$. 

This characteristic, which intensifies as MoE architectures deepen, results in significant fluctuations of the token-level importance ratios $w_{i,t}(\theta) = \frac{ \pi_{\theta} (y_{i,t} | x, y_{i,<t}) }{ \pi_{\theta_\text{old}} (y_{i,t} | x,y_{i,<t})}$. Such variability, as previously discussed in \S~\ref{sec:motivation} and~\ref{subsec:gradient}, renders these ratios unreliable and disrupts standard convergence during RL training.

\paragraph{Our Previous Approach} In our earlier work, we adopted the \textbf{Routing Replay} strategy as a training solution to this issue. This method involves storing the expert activations determined by $\pi_{\theta_\text{old}}$, then reusing these routing decisions within $\pi_{\theta}$ when calculating the importance ratios $w_{i,t}(\theta) = \frac{ \pi_{\theta} (y_{i,t} | x, y_{i,<t}) }{ \pi_{\theta_\text{old}} (y_{i,t} | x,y_{i,<t})}$. As a result, both $\pi_{\theta} (y_{i,t} | x, y_{i,<t})$ and $\pi_{\theta_\text{old}} (y_{i,t} | x,y_{i,<t})$ process $y_{i,t}$ using an identical expert configuration, thereby restoring stability to the token-wise importance ratios. This approach also preserves the consistency of the underlying network being optimized over successive gradient steps. As illustrated in Figure~\ref{fig:routing_replay}, Routing Replay plays a pivotal role in enabling reliable and stable convergence during GRPO training of MoE models.

\paragraph{Benefit of GSPO} While Routing Replay allows MoE models to achieve stable convergence during GRPO training, it does so at the cost of increased memory use and communication, and its reuse of routing modes restricts the effective capacity of the model. In contrast, as depicted in Figure~\ref{fig:results}, GSPO dispenses with the requirement for Routing Replay and is able to compute the importance ratios $s_i(\theta)$ through standard approaches, achieving consistent convergence and robust optimization. The central advantage lies in GSPO's focus on sequence likelihood (i.e., $\pi_{\theta} (y_i | x)$), making it less influenced by the likelihoods of individual tokens (i.e., $\pi_{\theta} (y_{i,t} | x, y_{i,<t})$). Given that MoE models inherently preserve their language modeling abilities, the sequence likelihood remains stable and does not exhibit extreme variation. Consequently, GSPO effectively addresses the challenge of unstable expert activation within MoE models, eliminating the necessity for intricate solutions such as Routing Replay. This results in a more straightforward and reliable training protocol, enabling models to take advantage of their complete capacity without imposing artificial limitations.

\subsection{Benefit of GSPO for RL Infrastructure}

Due to the varying numerical precisions present in training systems such as Megatron, compared to inference frameworks like SGLang and vLLM, it is common practice to rely on the training engine for recalculating the likelihoods of generated responses under the prior policy $\pi_{\theta_\text{old}}$. Nevertheless, GSPO relies on sequence-level likelihoods for optimization, rather than more granular token-level ones, which intuitively renders sequence-level computations less sensitive to differences in numerical precision. This property allows GSPO to utilize likelihood values produced by the inference engine directly within the optimization process, thereby eliminating the necessity to recalculate these values using the training engine. Such a feature proves particularly useful in contexts involving partial rollout, multi-turn reinforcement learning, and frameworks where training and inference occur on separate systems.

\section{Conclusion}

We introduce Group Sequence Policy Optimization (GSPO), a novel reinforcement learning approach designed for training extensive language models. GSPO applies importance sampling principles by calculating importance ratios grounded in sequence likelihoods, then implements sequence-level clipping, reward allocation, and optimization. Empirical results indicate GSPO achieves significantly greater stability, efficiency, and overall performance than GRPO, and proves especially effective for large-scale reinforcement learning applied to MoE models—serving as a key mechanism behind substantial improvements seen in recent Qwen3 models. Leveraging GSPO as a robust, scalable methodology, our future efforts will aim to further scale reinforcement learning, with the expectation of enabling fundamental progress in artificial intelligence.

\end{document}